\chapter{2004年全国IV卷（理）}
\section{单选题}
\begin{problem}
已知集合 \(M = \{0, 1, 2\}\)，\(N = \{x \mid x = 2a, a \in M\}\)，则集合 \(M \cap N =\)（\(\quad\)）

\choices
    {\( \{0\} \)}
    {\(\{0,1\}\)}
    {\(\{1,2\}\)}
    {\(\{0,2\}\)}
\end{problem}

\begin{answer}
D
\end{answer}

\begin{solution}
由 \(N=\{2a\mid a\in\{0,1,2\}\}=\{0,2,4\}\)，所以 \(M\cap N=\{0,2\}\)，故选 D.
\end{solution}

\begin{problem}
函数 \(y = \e^{2x}\)（\(x \in \R\)）的反函数为（\(\quad\)）

\choices
    {\(y = 2\ln x\)（\(x > 0\)）}
    {\(y = \ln\left(2x\right)\)（\(x > 0\)）}
    {\(y = \frac{1}{2}\ln x\)（\(x > 0\)）}
    {\(y = \frac{1}{2}\ln\left(2x\right)\)（\(x > 0\)）}
\end{problem}

\begin{answer}
C
\end{answer}

\begin{solution}
由 \(y=\e^{2x}\) 得 \(\ln y=2x\)，所以 \(x=\frac12\ln y\). 交换 \(x\)，\(y\) 后，反函数为 \(y=\frac12\ln x\)，其定义域为 \(x>0\)，故选 C.
\end{solution}

\begin{problem}
过点 \((-1, 3)\) 且垂直于直线 \(x - 2y + 3 = 0\) 的直线方程为（\(\quad\)）

\choices
    {\(2x + y - 1 = 0\)}
    {\(2x + y - 5 = 0\)}
    {\(x + 2y - 5 = 0\)}
    {\(x - 2y + 7 = 0\)}
\end{problem}

\begin{answer}
A
\end{answer}

\begin{solution}
直线 \(x-2y+3=0\) 的斜率为 \(\frac12\)，所以所求垂线的斜率为 \(-2\). 过 \((-1,3)\) 的直线方程为 \(y-3=-2(x+1)\)，整理得 \(2x+y-1=0\)，故选 A.
\end{solution}

\begin{problem}
\(\left(\frac{1 - \sqrt{3}\i}{1 + \i}\right)^2 =\)（\(\quad\)）

\choices
    {\(\sqrt{3} + \i\)}
    {\(-\sqrt{3} - \i\)}
    {\(\sqrt{3} - \i\)}
    {\(-\sqrt{3} + \i\)}
\end{problem}

\begin{answer}
D
\end{answer}

\begin{solution}
直接平方并利用 \((1+\i)^2=2\i\)，得
\[
\left(\frac{1-\sqrt{3}\i}{1+\i}\right)^2
=\frac{(1-\sqrt{3}\i)^2}{(1+\i)^2}
=\frac{-2-2\sqrt{3}\i}{2\i}
=-\sqrt{3}+\i.
\]
故选 D.
\end{solution}

\begin{problem}
不等式 \(\frac{x\left(x + 2\right)}{x - 3} < 0\) 的解集为（\(\quad\)）

\choices
    {\(\{x \mid x < -2 \text{或} 0 < x < 3\}\)}
    {\(\{x \mid -2 < x < 2 \text{或} x > 3\}\)}
    {\(\{x \mid x < -2 \text{或} x > 0\}\)}
    {\(\{x \mid x < 0 \text{或} x < 3\}\)}
\end{problem}

\begin{answer}
A
\end{answer}

\begin{solution}
临界点为 \(-2\)，\(0\) 和分母零点 \(3\). 在区间 \(\left(-\infty,-2\right)\)，\((-2,0)\)，\((0,3)\)，\(\left(3,+\infty\right)\) 上分别判断符号，分式的符号依次为负、正、负、正. 因为不等式是严格不等式，故解集为 \(\{x\mid x<-2\text{或}0<x<3\}\)，选 A.
\end{solution}

\begin{problem}
等差数列 \(\{a_{n}\}\) 中，\(a_{1} + a_{2} + a_{3} = -24\)，\(a_{18} + a_{19} + a_{20} = 78\)，则此数列前 \(20\) 项和等于（\(\quad\)）

\choices
    {\(160\)}
    {\(180\)}
    {\(200\)}
    {\(220\)}
\end{problem}

\begin{answer}
B
\end{answer}

\begin{solution}
等差数列中，三个连续项的和等于中间项的三倍，因此 \(3a_{2}=-24\)，得 \(a_{2}=-8\)，又 \(3a_{19}=78\)，得 \(a_{19}=26\). 于是公差 \(d=\frac{a_{19}-a_{2}}{17}=2\)，从而 \(a_{1}=-10\)，\(a_{20}=28\). 所以
\[
S_{20}=\frac{20(a_{1}+a_{20})}{2}=10\times18=180.
\]
故选 B.
\end{solution}

\begin{problem}
对于直线 \(m\)，\(n\) 和平面 \(\alpha\)，下面命题中的真命题是（\(\quad\)）

\choices
    {如果 \(m \subset \alpha\)，\(n \not\subset \alpha\)，\(m\)，\(n\) 是异面直线，那么 \(n \parallel \alpha\)}
    {如果 \(m \subset \alpha\)，\(n \not\subset \alpha\)，\(m\)，\(n\) 是异面直线，那么 \(n\) 与 \(\alpha\) 相交}
    {如果 \(m \subset \alpha\)，\(n \parallel \alpha\)，\(m\)，\(n\) 共面，那么 \(m \parallel n\)}
    {如果 \(m \parallel \alpha\)，\(n \parallel \alpha\)，\(m\)，\(n\) 共面，那么 \(m \parallel n\)}
\end{problem}

\begin{answer}
C
\end{answer}

\begin{solution}
在选项 C 中，\(n\parallel\alpha\)，所以 \(n\) 与 \(\alpha\) 没有公共点. 因为 \(m\subset\alpha\)，若共面的两直线 \(m\)，\(n\) 不平行，则它们必相交；交点同时属于 \(n\) 和 \(\alpha\)，矛盾. 因此 \(m\parallel n\)，选项 C 正确. 选项 A、B 中，异面直线与平面的关系既可能是平行，也可能相交；选项 D 中，两条都平行于同一平面的共面直线仍可能相交，故 A、B、D 均不成立.
\end{solution}

\begin{problem}
已知椭圆的中心在原点，离心率 \( e = \frac{1}{2} \)，且它的一个焦点与抛物线 \( y^{2} = -4x \) 的焦点重合，则此椭圆方程为（\(\quad\)）

\choices
    {\(\frac{x^2}{4} +\frac{y^2}{3} = 1\)}
    {\(\frac{x^2}{8} +\frac{y^2}{6} = 1\)}
    {\(\frac{x^2}{2} + y^2 = 1\)}
    {\(\frac{x^2}{4} + y^2 = 1\)}
\end{problem}

\begin{answer}
A
\end{answer}

\begin{solution}
抛物线 \(y^2=-4x\) 的焦点为 \((-1,0)\)，所以椭圆的焦距参数满足 \(c=1\). 由离心率 \(e=\frac ca=\frac12\)，得 \(a=2\). 再由 \(b^2=a^2-c^2\)，得 \(b^2=4-1=3\)，故椭圆方程为 \(\frac{x^2}{4}+\frac{y^2}{3}=1\)，选 A.
\end{solution}

\begin{problem}
从 \(5\) 位男教师和 \(4\) 位女教师中选出 \(3\) 位教师，派到 \(3\) 个班担任班主任（每班 \(1\) 位班主任），要求这 \(3\) 位班主任中男、女教师都要有，则不同的选派方案共有（\(\quad\)）

\choices
    {\(210\) 种}
    {\(420\) 种}
    {\(630\) 种}
    {\(840\) 种}
\end{problem}

\begin{answer}
B
\end{answer}

\begin{solution}
分两种情况计数. 若选出 \(2\) 位男教师和 \(1\) 位女教师，先选人再分配到班级，方案数为 \(\mathrm C_{5}^{2}\mathrm C_{4}^{1}\mathrm A_{3}^{3}\). 若选出 \(1\) 位男教师和 \(2\) 位女教师，方案数为 \(\mathrm C_{5}^{1}\mathrm C_{4}^{2}\mathrm A_{3}^{3}\). 因而总方案数为
\[
\mathrm C_{5}^{2}\mathrm C_{4}^{1}\mathrm A_{3}^{3}+\mathrm C_{5}^{1}\mathrm C_{4}^{2}\mathrm A_{3}^{3}
=10\times4\times6+5\times6\times6=420.
\]
故选 B.
\end{solution}

\begin{problem}
已知球的表面积为 \(20\pi\)，球面上有 \(A\)，\(B\)，\(C\) 三点. 如果 \(AB = AC = 2\)，\(BC = 2\sqrt{3}\)，则球心到平面 \(ABC\) 的距离为（\(\quad\)）

\choices
    {\(1\)}
    {\(\sqrt{2}\)}
    {\(\sqrt{3}\)}
    {\(2\)}
\end{problem}

\begin{answer}
A
\end{answer}

\begin{solution}
设球半径为 \(R\)，由 \(4\pi R^2=20\pi\)，得 \(R=\sqrt5\). 设三角形 \(ABC\) 外接圆半径为 \(r\). 由 \(AB=AC=2\)，\(BC=2\sqrt3\)，可知从 \(A\) 向 \(BC\) 作高，其高为 \(\sqrt{2^2-(\sqrt3)^2}=1\)，所以三角形面积为 \(\sqrt3\). 因此
\[
r=\frac{AB\cdot AC\cdot BC}{4[\triangle ABC]}=\frac{2\times2\times2\sqrt3}{4\sqrt3}=2.
\]
球心到平面 \(ABC\) 的垂线、球半径和该平面内的外接圆半径构成直角三角形. 设球心到平面的距离为 \(h\)，则 \(R^2=h^2+r^2\)，从而 \(h^2=5-4=1\)，即 \(h=1\)，选 A.
\end{solution}

\begin{problem}
\(\triangle ABC\) 中，\(a\)，\(b\)，\(c\) 分别为 \(\angle A\)，\(\angle B\)，\(\angle C\) 的对边. 如果 \(a\)，\(b\)，\(c\) 成等差数列，\(\angle B = 30^{\circ}\)，\(\triangle ABC\) 的面积为 \(\frac{3}{2}\)，那么 \(b =\)（\(\quad\)）

\choices
    {\(\frac{1 + \sqrt{3}}{2}\)}
    {\(1 + \sqrt{3}\)}
    {\(\frac{2 + \sqrt{3}}{2}\)}
    {\(2 + \sqrt{3}\)}
\end{problem}

\begin{answer}
B
\end{answer}

\begin{solution}
因为 \(a\)，\(b\)，\(c\) 成等差数列，所以 \(a+c=2b\). 由面积公式
\[
\frac12ac\sin B=\frac32
\]
及 \(\sin30^{\circ}=\frac12\)，得 \(ac=6\). 再由余弦定理，
\[
b^2=a^2+c^2-2ac\cos30^{\circ}
=(a+c)^2-(2+\sqrt3)ac
=4b^2-6(2+\sqrt3).
\]
所以 \(b^2=2(2+\sqrt3)=4+2\sqrt3=(1+\sqrt3)^2\). 因为边长为正，故 \(b=1+\sqrt3\)，选 B.
\end{solution}

\begin{problem}
设函数 \(f(x)\)（\(x \in \R\)）为奇函数，\(f(1) = \frac{1}{2}\)，\(f(x + 2) = f(x) + f(2)\)，则 \(f(5) =\)（\(\quad\)）

\choices
    {\(0\)}
    {\(1\)}
    {\( \frac{5}{2} \)}
    {\(5\)}
\end{problem}

\begin{answer}
C
\end{answer}

\begin{solution}
因 \(f\) 为奇函数，\(f(0)=0\)，且在题设关系中令 \(x=-1\)，得 \(f(1)=f(-1)+f(2)=-f(1)+f(2)\). 因而 \(f(2)=2f(1)=1\). 再由递推关系，\(f(3)=f(1)+f(2)=\frac32\)，\(f(5)=f(3)+f(2)=\frac52\)，故选 C.
\end{solution}

\section{填空题}
\begin{problem}
\(\left(x - \frac{1}{\sqrt{x}}\right)^8\) 展开式中 \(x^5\) 的系数为\fillinblank{}.
\end{problem}

\begin{answer}
\(28\)
\end{answer}

\begin{solution}
展开式的通项为
\[
\mathrm C_{8}^{k}x^{8-k}\left(-\frac1{\sqrt{x}}\right)^k
=(-1)^k\mathrm C_{8}^{k}x^{8-\frac32k}.
\]
令指数满足 \(8-\frac32k=5\)，得 \(k=2\). 因此所求系数为 \((-1)^2\mathrm C_{8}^{2}=28\).
\end{solution}

\begin{problem}
向量 \(\bs{a}\)，\(\bs{b}\) 满足 \(\left(\bs{a} - \bs{b}\right) \cdot \left(2\bs{a} + \bs{b}\right) = -4\)，且 \(|\bs{a}| = 2\)，\(|\bs{b}| = 4\)，则 \(\bs{a}\) 与 \(\bs{b}\) 夹角的余弦值等于\fillinblank{}.
\end{problem}

\begin{answer}
\(-\frac12\)
\end{answer}

\begin{solution}
展开数量积，得
\[
(\bs{a}-\bs{b})\cdot(2\bs{a}+\bs{b})
=2|\bs{a}|^2-\bs{a}\cdot\bs{b}-|\bs{b}|^2=-4.
\]
代入 \(|\bs{a}|=2\)，\(|\bs{b}|=4\)，得到 \(8-\bs{a}\cdot\bs{b}-16=-4\)，所以 \(\bs{a}\cdot\bs{b}=-4\). 设夹角为 \(\theta\)，则
\[
\cos\theta=\frac{\bs{a}\cdot\bs{b}}{|\bs{a}|\,|\bs{b}|}=\frac{-4}{2\times4}=-\frac12.
\]
\end{solution}

\begin{problem}
函数 \(f(x) = \cos x - \frac{1}{2}\cos 2x\)（\(x \in \R\)）的最大值等于\fillinblank{}.
\end{problem}

\begin{answer}
\(\frac34\)
\end{answer}

\begin{solution}
令 \(t=\cos x\)，则 \(-1\leqslant t\leqslant1\)，且 \(\cos2x=2t^2-1\). 所以
\[
f(x)=t-\frac12(2t^2-1)=\frac34-\left(t-\frac12\right)^2\leqslant\frac34.
\]
当 \(t=\frac12\) 时等号成立，故最大值为 \(\frac34\).
\end{solution}

\begin{problem}
设 \(x\)，\(y\) 满足约束条件：\(\begin{cases}
x + y \leqslant 1,\\
y \leqslant x,\\
y \geqslant 0.
\end{cases}\) 则 \(z = 2x + y\) 的最大值是\fillinblank{}.
\end{problem}

\begin{answer}
\(2\)
\end{answer}

\begin{solution}
可行域由直线 \(x+y=1\)，\(y=x\) 和 \(y=0\) 围成，三个顶点分别为 \((0,0)\)，\((1,0)\) 和 \(\left(\frac12,\frac12\right)\). 在这些顶点处，\(z=2x+y\) 的值分别为 \(0\)，\(2\) 和 \(\frac32\). 因此 \(z\) 的最大值为 \(2\).
\end{solution}

\section{解答题}
\begin{problem}
已知 \(\alpha\) 为第二象限角，且 \(\sin \alpha = \frac{\sqrt{15}}{4}\)，求 \(\frac{\sin\left(\alpha + \frac{\pi}{4}\right)}{\sin 2\alpha + \cos 2\alpha + 1}\) 的值.
\end{problem}

\begin{solution}
因为 \(\alpha\) 为第二象限角，且 \(\sin\alpha=\frac{\sqrt{15}}4\)，所以 \(\cos\alpha=-\frac14\). 从而
\[
\sin\left(\alpha+\frac\pi4\right)=\frac{\sin\alpha+\cos\alpha}{\sqrt2}=\frac{\sqrt{15}-1}{4\sqrt2},
\]
并且 \(\sin2\alpha=2\sin\alpha\cos\alpha=-\frac{\sqrt{15}}8\)，\(\cos2\alpha=\cos^2\alpha-\sin^2\alpha=-\frac78\). 因此
\[
\sin2\alpha+\cos2\alpha+1=\frac{1-\sqrt{15}}8.
\]
所求值为
\[
\frac{\dfrac{\sqrt{15}-1}{4\sqrt2}}{\dfrac{1-\sqrt{15}}8}=-\sqrt2.
\]
\end{solution}

\begin{problem}
求函数 \(f(x) = \ln\left(1 + x\right) - \frac{1}{4} x^2 \) 在 \([0,2]\) 上的最大值和最小值.
\end{problem}

\begin{solution}
函数在 \([0,2]\) 上可导，且
\[
f'(x)=\frac1{1+x}-\frac{x}{2}=\frac{2-x-x^2}{2(1+x)}=-\frac{(x+2)(x-1)}{2(x+1)}.
\]
由于 \(x+1>0\)，所以 \(f'(x)>0\) 在 \([0,1)\) 上成立，\(f'(x)<0\) 在 \((1,2]\) 上成立. 因此函数在 \(x=1\) 处取得最大值
\[
f(1)=\ln2-\frac14.
\]
再比较端点值，\(f(0)=0\)，\(f(2)=\ln3-1>0\)，所以最小值为 \(0\). 故最大值为 \(\ln2-\frac14\)，最小值为 \(0\).
\end{solution}

\begin{problem}
某同学参加科普知识竞赛，需回答 \(3\) 个问题. 竞赛规则规定：每题回答正确得 \(100\) 分，回答不正确得 \(-100\) 分. 假设这名同学每题回答正确的概率均为 \(0.8\)，且各题回答正确与否相互之间没有影响.

\begin{enumerate}
\item 求这名同学回答这三个问题的总得分 \(\xi\) 的概率分布和数学期望；
\item 求这名同学总得分不为负分（即 \(\xi \geqslant 0\)）的概率.
\end{enumerate}
\end{problem}

\begin{solution}
设答对题数为 \(K\)，则 \(K\sim B\left(3,\frac45\right)\)，且总得分为 \(\xi=100K-100(3-K)=200K-300\). 因此
\begin{center}
\begin{tabular}{c|cccc}
\(\xi\) & \(-300\) & \(-100\) & \(100\) & \(300\) \\ \hline
\(P\) & \(\dfrac1{125}\) & \(\dfrac{12}{125}\) & \(\dfrac{48}{125}\) & \(\dfrac{64}{125}\)
\end{tabular}
\end{center}
其中各概率分别由 \(P(K=k)=\mathrm C_{3}^{k}\left(\frac45\right)^k\left(\frac15\right)^{3-k}\) 得到. 数学期望为
\[
E(\xi)=(-300)\frac1{125}+(-100)\frac{12}{125}+100\frac{48}{125}+300\frac{64}{125}=180.
\]
当 \(\xi\geqslant0\) 时，必须有 \(K\geqslant2\)，故
\[
P(\xi\geqslant0)=P(K=2)+P(K=3)=\frac{48+64}{125}=\frac{112}{125}.
\]
\end{solution}

\begin{problem}
如图，四棱锥 \(P - ABCD\) 中，底面 \(ABCD\) 为矩形，\(AB = 8\)，\(AD = 4\sqrt{3}\)，侧面 \(PAD\) 为等边三角形，并且与底面所成二面角为 \(60^{\circ}\).

\begin{enumerate}
\item 求四棱锥 \(P - ABCD\) 的体积；
\item 证明：\(PA \perp BD\).
\end{enumerate}

\begin{center}
\bitmapfigure[max width=0.42\linewidth]{img/2004/national_paper_4_science/q20_fig1.png}
\end{center}
\end{problem}

\begin{solution}
设 \(M\) 为 \(AD\) 的中点. 因为 \(\triangle PAD\) 为边长 \(4\sqrt3\) 的等边三角形，所以 \(PM\perp AD\)，且 \(PM=\frac{\sqrt3}{2}AD=6\). 过 \(PM\) 作垂直于 \(AD\) 的截面，截得的两平面夹角为 \(60^{\circ}\). 因而点 \(P\) 到底面的距离为
\[
h=PM\sin60^{\circ}=6\times\frac{\sqrt3}{2}=3\sqrt3.
\]
底面面积为 \(8\times4\sqrt3=32\sqrt3\)，所以四棱锥体积为
\[
V=\frac13\times32\sqrt3\times3\sqrt3=96.
\]

建立直角坐标系，使 \(A=(0,0,0)\)，\(B=(8,0,0)\)，\(D=(0,4\sqrt3,0)\)，且 \(z\) 轴垂直于底面. 由上面的截面关系，点 \(P\) 在底面的射影为 \(H=(3,2\sqrt3,0)\)，故 \(P=(3,2\sqrt3,3\sqrt3)\). 于是
\[
\overrightarrow{AP}=(3,2\sqrt3,3\sqrt3),\qquad \overrightarrow{BD}=(-8,4\sqrt3,0),
\]
从而
\[
\overrightarrow{AP}\cdot\overrightarrow{BD}=3\times(-8)+2\sqrt3\times4\sqrt3=0.
\]
因此 \(AP\perp BD\)，即 \(PA\perp BD\).
\end{solution}

\begin{problem}
双曲线 \(\frac{x^2}{a^2} - \frac{y^2}{b^2} = 1\)（\(a > 1\)，\(b > 0\)）的焦点距为 \(2c\)，直线 \(l\) 过点 \((a,0)\) 和 \((0,b)\)，且点 \((1,0)\) 到直线 \(l\) 的距离与点 \((-1,0)\) 到直线 \(l\) 的距离之和 \(s \geqslant \frac{4}{5} c\). 求双曲线的离心率 \(e\) 的取值范围.
\end{problem}

\begin{solution}
直线 \(l\) 的方程为 \(\frac{x}{a}+\frac{y}{b}=1\)，即 \(bx+ay-ab=0\). 由于 \(a>1\)，两点到直线的距离之和为
\[
s=\frac{b(a-1)}{\sqrt{a^2+b^2}}+\frac{b(a+1)}{\sqrt{a^2+b^2}}=\frac{2ab}{c},
\]
其中双曲线满足 \(c^2=a^2+b^2\). 代入题设不等式，得
\[
\frac{2ab}{c}\geqslant\frac45c
\Longleftrightarrow 5ab\geqslant2(a^2+b^2).
\]
令 \(t=\frac ba>0\)，约去 \(a^2\) 后得到
\[
2t^2-5t+2\leqslant0
\Longleftrightarrow (2t-1)(t-2)\leqslant0,
\]
所以 \(\frac12\leqslant t\leqslant2\). 双曲线离心率为
\[
e=\frac ca=\sqrt{1+\frac{b^2}{a^2}}=\sqrt{1+t^2}.
\]
由于 \(t>0\)，该函数在所给区间上递增，故
\[
\frac{\sqrt5}{2}=\sqrt{1+\frac14}\leqslant e\leqslant\sqrt5.
\]
\end{solution}

\begin{problem}
已知函数 \(f(x) = \e^{-x}\left(\cos x + \sin x\right)\). 将满足 \(f'(x) = 0\) 的所有正数 \(x\) 从小到大排成数列 \(\{x_n\}\).

\begin{enumerate}
\item 证明数列 \(\{f(x_n)\}\) 为等比数列；
\item 记 \(S_{n}\) 是数列 \(\{x_n f(x_n)\}\) 的前 \(n\) 项和，求 \(\displaystyle\lim_{n\to \infty}\frac{S_1 + S_2 + \cdots + S_n}{n}\).
\end{enumerate}
\end{problem}

\begin{solution}
\begin{enumerate}
\item 对函数求导，得
\[
f'(x)=-\e^{-x}(\cos x+\sin x)+\e^{-x}(-\sin x+\cos x)=-2\e^{-x}\sin x.
\]
因为 \(\e^{-x}\ne0\)，所以正数零点满足 \(\sin x=0\)，从而 \(x_n=n\pi\)（\(n=1,2,\ldots\)）. 因此
\[
f(x_n)=\e^{-n\pi}\cos(n\pi)=(-1)^n\e^{-n\pi},
\]
且 \(\frac{f(x_{n+1})}{f(x_n)}=-\e^{-\pi}\) 为常数，所以 \(\{f(x_n)\}\) 是等比数列，首项为 \(-\e^{-\pi}\)，公比为 \(-\e^{-\pi}\).

\item 由 \(x_n=n\pi\)，有
\[
S_n=\sum_{k=1}^{n}x_kf(x_k)=\pi\sum_{k=1}^{n}k(-\e^{-\pi})^k.
\]
由于 \(\sum_{k=1}^{\infty}k\e^{-k\pi}\) 收敛，故 \(S_n\) 收敛，且
\[
\lim_{n\to\infty}S_n
=\pi\sum_{k=1}^{\infty}k(-\e^{-\pi})^k
=\pi\frac{-\e^{-\pi}}{\left(1+\e^{-\pi}\right)^2}.
\]
数列 \(S_n\) 收敛时，其算术平均数也收敛到同一极限，因此
\[
\lim_{n\to\infty}\frac{S_1+S_2+\cdots+S_n}{n}
=-\frac{\pi\e^{-\pi}}{\left(1+\e^{-\pi}\right)^2}.
\]
\end{enumerate}
\end{solution}
